\section{Unternehmen}

TrainCloud GmbH bietet Software-as-a-Services-(SaaS)-Lösungen für die Eisenbahnindustrie an. 
Das Unternehmen entwickelt digitale Lösungen zur Unterstützung der Verwaltung, Überwachung 
und Durchführung von Betriebsprozessen im Eisenbahnbereich. 

Eines der aktueller Produkte des Unternehmens ist eine interaktive, cloudbasierte Plattform 
zur Überwachung des Eisenbahnverkehrs. Sie ermöglicht es Kunden, betriebsbezogene Informationen 
zentral zu verfolgen und zu verwalten. 

Zudem entwickelt TrainCloud eine mobile Anwendung für Eisenbahnbetriebspersonal wie Triebfahrzeugführer, 
Wagenmeister und Rangierbegleiter. Die Anwendung unterstützt das Personal bei der Durchführung und Verwaltung 
der täglichen betrieblichen Aufgaben, beispielsweise bei der Erstellung und Verwaltung verschiedener Meldungen 
und Berichte, darunter Verspätungs-, Schadenmelsungen sowie weiterer damit verbundener betrieblicher Prozesse.

Im Eisenbahnbereich werden einige Betriebs- und Verwaltungsprozesse nach wie vor auf traditionelle, papierbasierte Weise abgewickelt. 
Dies kann den Verwaltungsaufwand erhöhen und die Speicherung, Verarbeitung und zentrale Verwaltung von Daten erschweren. 

Vor dem Hintergrund der zunehmenden Digitalisierung gewinnt die Konvertierung traditioneller Methoden in digitale Lösungen 
immer mehr an Bedeutung. TrainCloud verfolgt das Ziel, traditionelle Verwaltungs- und Meldeprozesse im Eisenbahnbereich zu digitalisieren. 
Dabei arbeitet das Entwicklungsteam eng mit dem Eisenbahnbetriebspersonal zusammen. 
Durch diese Zusammenarbeit sollen Softwarelösungen entwickelt werden, die den Anforderungen 
und den tatsächlichen Arbeitsabläufen im Eisenbahnbetrieb entsprechen.

TrainCloud bietet zwei sich ergänzende digitale Produkte an:
\begin{itemize}
    \item \textbf{TrainCloud Web Portal:} Eine cloudbasierte Verwaltungsplattform für Eisenbahnunternehmen.
    Sie bietet Führungskräften und Administratoren einen zentralen Überblick über die betrieblichen Daten des Unternehmens. 
    Dazu gehören unter anderem zusammengefasste Berichte, Personalverwaltung und Fahrzeugverwaltung. 

    \item \textbf{TrainCloud Mobile Application:} Eine mobile Anwendung für das einzelne Eisenbahbetriebspersonal, darunter Triebfahrzeugführer, 
    Wagenmeister und Rangierbegleiter. Die Anwendung unterstützt die Mitarbeiter bei der Verwaltung und Dutchführung ihrer taglichen betrieblichen 
    Aufgaben.
\end{itemize}

Zusammen bilden diese beiden Produkte eine integriertes digites System, das das im operativen Bereich tägtige Personal mit der Unternehmensverwaltung
auf einer gemeinsamen Plattform verbindet.

\section{Praktikumstelle}

\section{Programmierungsprache}

\section{Zusammenfassung der fertigen gemachten Projekten}